\documentclass[11pt,article,oneside]{memoir}
\usepackage[utf8]{inputenc}
\usepackage[T1]{fontenc}
\usepackage{url}
\usepackage[hidelinks,colorlinks=true,linkcolor=blue]{hyperref}
\usepackage{amsmath,amssymb}
\usepackage{graphicx}
\usepackage{xcolor}
\usepackage{tabularx}
\usepackage{listings}
\usepackage{geometry}
\usepackage{titlesec}
\usepackage{libertine}

\geometry{margin=1in}

\titleformat{\section}{\Large\bfseries\color{darkgray}}{\thesection}{1em}{}
\titleformat{\subsection}{\large\bfseries\color{blue!80!black}}{\thesubsection}{1em}{}

\lstset{
  basicstyle=\small\ttfamily,
  breaklines=true,
  frame=lines,
  backgroundcolor=\color{gray!10}
}

\begin{document}

\begin{center}
    \Huge \textbf{Reasoning Studio Guide} \\
    \vspace{0.5em}
    \large \textit{A CEREBRUM Research Publication (v1.2.4)} \\
    \vspace{1em}
    \large \textbf{Bryan Alexander Buchorn} \\
    Independent Researcher \\
    \vspace{2em}
\end{center}

\begin{abstract}
\textbf{Overview:} This guide provides a conceptual entry point into the CEREBRUM framework. It is designed for practitioners and researchers seeking to understand the core reasoning signals of Framework Implementation Guide.
\end{abstract}

\section{CEREBRUM Reasoning Studio — User Guide}

\textbf{Version}: v1.1.0
\textbf{Interface}: Gradio web UI + pyvis graph visua\-lization
\textbf{Module}: \texttt{ui/studio.py}

---

\subsection{Overview}

The Reasoning Studio is CEREBRUM's interactive visual interface. It lets you load a Knowledge Graph, run multi-hop reasoning queries, inspect every step of the CSA attention mechanism, and watch community structure form in real time — all without writing code.

This is the "Glass-Box" in action: you can see exactly what the AI is doing at every reasoning step.

---

\subsection{Instal\-lation}

\begin{lstlisting}
pip install -r ui/requirements.txt
\end{lstlisting}

Depend\-encies installed:
\begin{itemize}
\item \texttt{gradio} — web UI framework
\item \texttt{pyvis} — interactive graph visua\-lization
\item \texttt{networkx} — graph backend (already in core)

\end{itemize}
---

\subsection{Starting the Studio}

\begin{lstlisting}
python ui/studio.py
\end{lstlisting}

The Studio opens at \texttt{http://localhost:7860} in your browser. No API server is required — the Studio embeds the full CEREBRUM engine directly.

\subsubsection{Command-line options}

\begin{lstlisting}
# Specify a graph file to load on startup
python ui/studio.py --csv path/to/graph.csv

# Use semantic embeddings (requires sentence-transformers)
python ui/studio.py --embeddings sentence

# Change port
python ui/studio.py --port 8080

# Share publicly (Gradio tunnel — for demos)
python ui/studio.py --share
\end{lstlisting}

---

\subsection{Interface Overview}

The Studio has four main panels:

\subsubsection{1. Graph Panel (left)}
An interactive force-directed visua\-lization of the loaded Knowledge Graph powered by pyvis.

\begin{itemize}
\item \textbf{Nodes} are colored by community (each community gets a distinct color)
\item \textbf{Edges} show relation types as labels
\item \textbf{Node size} reflects PageRank centrality — important nodes appear larger
\item \textbf{Click any node} to set it as the query start entity
\item \textbf{Hover} over an edge to see its CSA attention weight

\end{itemize}
\subsubsection{2. Query Panel (top center)}
Controls for running reasoning queries:

\begin{center}
\begin{tabularx}{\columnwidth}{|>{\raggedright\arraybackslash}X|>{\raggedright\arraybackslash}X|}
\hline
Control & Description \\ \hline
\textbf{Entity} & Starting node for traversal \\ \hline
\textbf{Max Hops} & Maximum reasoning depth (1–5) \\ \hline
\textbf{Beam Width} & Candidates per hop (1–50) \\ \hline
\textbf{Top-K} & Number of answer paths to return \\ \hline
\textbf{Probab\-ilistic} & Enable Bayesian beam search \\ \hline
\textbf{Warm-Start} & Beta prior seeding strength (0–10) \\ \hline
\end{tabularx}
\end{center}

Click \textbf{Run Query} to execute. Results appear in the Answer Panel within milli\-seconds.

\subsubsection{3. Answer Panel (center)}
Displays ranked answer paths with:
\begin{itemize}
\item \textbf{Entity}: the terminal node of each path
\item \textbf{Score}: composite CSA path score (0–1)
\item \textbf{Path}: the full hop-by-hop trace (\texttt{A → rel → B → rel → C})
\item \textbf{Confidence interval}: when proba\-bilistic mode is on
\item \textbf{Community trace}: which communities were traversed at each hop

\end{itemize}
Click any answer row to \textbf{highlight the reasoning path} in the Graph Panel.

\subsubsection{4. Attention Math Panel (right)}
When you click an edge in the Graph Panel or an answer path in the Answer Panel, this panel shows the full CSA formula breakdown for that specific edge:

\begin{lstlisting}
a(Marie Curie → Radium, hop=1) = sigmoid(
  0.4 × cos(e_u, e_v) = 0.4 × 0.71  =  0.284
  0.4 × S_C(u,v)      = 0.4 × 1.00  =  0.400   ← same community
  0.1 × w_rel         = 0.1 × 0.95  =  0.095   ← "discovered" weight
  0.05 × d_norm       = 0.05 × 0.10 = -0.005
  0.05 × φ(k)         = 0.05 × 1.00 =  0.050   ← hop 1
  0.1 × PR(v)         = 0.1 × 0.43  =  0.043
                        ─────────────────────
                        total input           =  0.867
  sigmoid(0.867)                              =  0.704
)
\end{lstlisting}

This is the "Forensic Math Panel" — you see exactly why this edge received its score.

---

\subsection{Loading a Graph}

\subsubsection{From a CSV file}
Click \textbf{Upload CSV} in the Graph Panel toolbar. Format: \texttt{subject,predicate,object[,weight]}

\subsubsection{From the toy graph (built-in)}
Click \textbf{Load Toy Graph} to load the canonical 21-node test graph from \texttt{tests/fixtures/toy\_graph.csv}. Useful for exploring the interface before loading your own data.

\subsubsection{From a URL (Remote CSV)}
Enter a URL in the \textbf{Graph URL} field and click \textbf{Load}. The Studio fetches and parses the CSV remotely.

---

\subsection{Community Visual\-ization}

Click the \textbf{Communities} tab to switch the Graph Panel to community-focus mode:
\begin{itemize}
\item Each community is shown as a convex hull overlay
\item Community sizes and modularity Q score are displayed in the sidebar
\item Use the \textbf{Rebalance} button to trigger a fresh DSCF run and watch communities reorganize

\end{itemize}
---

\subsection{Interactive Attention Tuning (Dialectics)}

The \textbf{Attention Sliders} section lets you adjust the six CSA weights ($\alpha$, $\beta$, $\gamma$, $\delta$, $\epsilon$, $\zeta$) in real time using sliders. When you move a slider, the last query autom\-atically re-runs with the new weights and the Graph Panel updates to show how the reasoning path changes.

This allows you to answer questions like:
\begin{itemize}
\item "What if I trust community structure more than semantic similarity?" (increase $\beta$)
\item "What if I weight relation types more heavily?" (increase $\gamma$)
\item "What if I penalize long paths more?" (increase $\delta$)

\end{itemize}
---

\subsection{Live Streaming Mode}

Click the \textbf{Stream} tab to enable streaming ingest mode:
\begin{enumerate}
\item Select a discretizer (Threshold, STDP, Delta, Frequency, Pattern)
\item Configure thresholds
\item Click \textbf{Start Stream}
\item Watch edges materialize in the Graph Panel in real time as events are processed
\item Community structure auto-updates when modularity drift triggers a rebalance

\end{enumerate}
---

\subsection{Exporting Results}

\begin{itemize}
\item \textbf{Export Path} — saves the selected reasoning trace as a JSON file including the full CSA math breakdown
\item \textbf{Export Graph} — saves the current graph (with community assignments) as a CSV
\item \textbf{Screenshot} — captures the current Graph Panel view as a PNG

\end{itemize}
---

\subsection{Keyboard Shortcuts}

\begin{center}
\begin{tabularx}{\columnwidth}{|>{\raggedright\arraybackslash}X|>{\raggedright\arraybackslash}X|}
\hline
Key & Action \\ \hline
\texttt{Enter} & Run query with current settings \\ \hline
\texttt{Esc} & Clear selection \\ \hline
\texttt{Space} & Toggle community overlay \\ \hline
\texttt{R} & Rebalance communities \\ \hline
\texttt{→} & Step through reasoning path (hop by hop) \\ \hline
\end{tabularx}
\end{center}

---

\subsection{Troubl\-eshooting}

\textbf{Studio won't start}: Ensure \texttt{gradio} and \texttt{pyvis} are installed: \texttt{pip install -r ui/requirements.txt}

\textbf{Graph not rendering}: Large graphs (\textgreater 5,000 nodes) may take several seconds for pyvis to lay out. Use the \textbf{Limit Display} slider to show only the top-N nodes by PageRank.

\textbf{No paths found}: The query entity may not exist in the loaded graph. Try the autoc\-omplete dropdown or check entity names are exact matches (case-sensitive).

\textbf{Slow queries}: Reduce \texttt{beam\_width} or \texttt{max\_hops} in the Query Panel.

---
\textbf{Copyright © 2026 Bryan Alexander Buchorn. All Rights Reserved.}


\vspace{3em}
\begin{center}
    \small \textit{Project CEREBRUM — March 2026 — Hardened Enterprise Security (v1.2.4)}
\end{center}

\end{document}
